\documentclass[11pt,a4paper]{article}
% preamble.tex — shared setup for human-verification documents
% Compile each document with:  tectonic <File>.tex

\usepackage[T1]{fontenc}
\usepackage[utf8]{inputenc}
\usepackage{lmodern}
\usepackage{amssymb,amsmath,amsthm}
\usepackage{mathtools}
\usepackage[dvipsnames]{xcolor}
\usepackage{listings}
\usepackage{hyperref}
\usepackage[margin=2.5cm]{geometry}
\usepackage{parskip}

% ── Colours ──────────────────────────────────────────────────────────
\definecolor{leanBg}{HTML}{F7F7F0}
\definecolor{haskellBg}{HTML}{F0F4F7}
\definecolor{keyword}{HTML}{7B3294}
\definecolor{comment}{HTML}{5A8C32}
\definecolor{noteFg}{HTML}{B35900}

% ── Lean 4 listings language ─────────────────────────────────────────
% NOTE: Lean listing blocks use ASCII approximations of unicode symbols
% (e.g. -> for arrow, o. for compose).  The literate key substitutes
% these with proper math-mode renderings in the PDF.
\lstdefinelanguage{Lean4}{
  morekeywords={def,theorem,structure,where,namespace,end,import,fun,
                let,have,show,by,exact,intro,cases,rfl,subst,obtain,
                funext,simp,rw,apply,variable,Type,Prop,instance,
                section,open,class,extends,deriving,sorry,noncomputable},
  sensitive=true,
  morecomment=[l]{--},
  morecomment=[n]{/-}{-/},
  morestring=[b]",
}

% ── Lean listing style ───────────────────────────────────────────────
\lstdefinestyle{lean}{
  language=Lean4,
  backgroundcolor=\color{leanBg},
  basicstyle=\small\ttfamily,
  keywordstyle=\color{keyword}\bfseries,
  commentstyle=\color{comment}\itshape,
  breaklines=true,
  frame=single,
  framerule=0.4pt,
  xleftmargin=4pt,
  xrightmargin=4pt,
  aboveskip=6pt,
  belowskip=4pt,
  literate=
    {->}{{$\to$}}2
    {=>}{{$\Rightarrow$}}2
    {<|}{{$\langle$}}1
    {|>}{{$\rangle$}}1
    {==>}{{$\Longrightarrow$}}3
    { o. }{{$\;\circ\;$}}3
    {forall}{{$\forall$}}1
    {:=}{{$\,:=\,$}}2
    {!=}{{$\neq$}}2
    {<=}{{$\leq$}}2
    {>=}{{$\geq$}}2,
}

% ── Haskell listing style ───────────────────────────────────────────
\lstdefinestyle{haskell}{
  language=Haskell,
  backgroundcolor=\color{haskellBg},
  basicstyle=\small\ttfamily,
  keywordstyle=\color{keyword}\bfseries,
  commentstyle=\color{comment}\itshape,
  breaklines=true,
  frame=single,
  framerule=0.4pt,
  xleftmargin=4pt,
  xrightmargin=4pt,
  aboveskip=6pt,
  belowskip=4pt,
  literate=
    {->}{{$\to$}}2
    {=>}{{$\Rightarrow$}}2,
}

% ── Theorem environments ────────────────────────────────────────────
\theoremstyle{definition}
\newtheorem{defn}{Definition}[section]
\newtheorem{thm}[defn]{Theorem}
\newtheorem{prop}[defn]{Proposition}

\theoremstyle{remark}
\newtheorem*{fidelitynote}{Fidelity note}

% ── Convenience macros ───────────────────────────────────────────────
\newcommand{\lean}[1]{\texttt{#1}}      % inline Lean identifier
\newcommand{\hask}[1]{\texttt{#1}}      % inline Haskell identifier
\newcommand{\Ext}{\mathrm{Ext}}
\newcommand{\Pos}{\mathrm{Pos}}
\newcommand{\Shape}{\mathrm{Shape}}


\title{Human Verification: \texttt{Containers.ProofSearchN}\\[4pt]
  \large N-ary Proof Search Trees as a Directed Container}
\author{Machine-generated from \texttt{lean/ProofSearchN.lean} for review by Neil Ghani}
\date{July 2026}

\begin{document}
\maketitle

\begin{abstract}
This document presents every definition and theorem \emph{statement} from the
Lean~4 module \lean{Containers.ProofSearchN} in three parallel forms: the
original Lean, a readable Haskell rendering, and precise mathematical English.
Proofs are omitted --- they are machine-checked by Lean.  The purpose is to
let a human verifier confirm that the formalisation captures the intended
mathematics.
\end{abstract}

\tableofcontents
\bigskip

%===================================================================
\section{Overview}
%===================================================================

This file generalises \lean{Containers.ProofSearch} from \emph{binary} trees
to \textbf{finitely-branching rose trees} (arbitrary arity at each node).  The
directed-container structure is identical in spirit: shapes are rose trees,
positions are paths from root to a node, and the comonad operations come from
path concatenation.

\paragraph{Encoding choice.}
Each node carries a branching factor $n : \mathbb{N}$ and a function assigning
a child tree to each index $i : \mathrm{Fin}(n)$.  The \lean{Fin\,n} encoding
avoids nested-inductive complications (which \lean{List} would introduce) and
lets structural recursion work directly on \lean{RTree} and \lean{RPos} without
well-founded recursion or fuel arguments.

A \textbf{leaf} (dead end or success) is \lean{node\,0\,f} where
$f : \mathrm{Fin}(0) \to \mathrm{RTree}$ is vacuously defined.  The only
position in a leaf is \lean{here}.

\paragraph{Relationship to \lean{STree}.}
\lean{STree} embeds into \lean{RTree}: a leaf maps to \lean{node\,0\,!} and
\lean{node\,$\ell$\,$r$} maps to \lean{node\,2\,![$\ell'$, $r'$]}.  This is
not a strict sub-case since \lean{STree} forces exactly 0 or 2 children.

%===================================================================
\section{Shapes: finitely-branching rose trees}
%===================================================================

% --- 1. RTree -----------------------------------------------------
\subsection{Inductive type: \lean{RTree}}

\begin{defn}[Rose tree]
A \textbf{finitely-branching rose tree} is a node with a branching factor
$n : \mathbb{N}$ and a family of $n$ child trees:
\[
  \mathrm{RTree} \;::=\; \lean{node}(n : \mathbb{N},\; \mathrm{Fin}(n) \to \mathrm{RTree}).
\]
A leaf (dead end) is $\lean{node}(0, f)$ where $f$ is vacuous.
\end{defn}

\paragraph{Lean.}
\begin{lstlisting}[style=lean]
inductive RTree where
  | node : (n : Nat) -> (Fin n -> RTree) -> RTree
\end{lstlisting}

\paragraph{Haskell.}
\begin{lstlisting}[style=haskell]
data RTree = Node Int [RTree]   -- n children stored as a list
  deriving (Show, Eq)

-- A leaf is Node 0 [].
-- In Lean, children are indexed by Fin n (a bounded integer type).
-- Haskell approximates this with a plain list.
\end{lstlisting}

\begin{fidelitynote}
In Lean, the children are given by a function $\mathrm{Fin}(n) \to \mathrm{RTree}$
where $n$ is stored in the tree.  This allows structural recursion on
\lean{RTree} and \lean{RPos} directly.  In Haskell, a list of children is
the natural representation but cannot enforce the branching-factor invariant
at the type level.
\end{fidelitynote}


%===================================================================
\section{Positions: paths in a rose tree}
%===================================================================

% --- 2. RPos ------------------------------------------------------
\subsection{Inductive type: \lean{RPos}}

\begin{defn}[Position in a rose tree]
A \textbf{position} in an \lean{RTree} is a path from the root to a node:
\begin{itemize}
  \item \lean{here}: the root (valid for any tree);
  \item \lean{child}$(i, p)$: descend into child $i$ (where $i : \mathrm{Fin}(n)$)
        and continue along path $p$ in that child.
\end{itemize}
For a leaf (\lean{node\,0\,f}), there are no valid $i : \mathrm{Fin}(0)$,
so the only position is \lean{here}.
\end{defn}

\paragraph{Lean.}
\begin{lstlisting}[style=lean]
inductive RPos : RTree -> Type where
  | here  : RPos t
  | child : (i : Fin n) -> RPos (children i) -> RPos (.node n children)
\end{lstlisting}

\paragraph{Haskell.}
\begin{lstlisting}[style=haskell]
data RPos t where
  Here  :: RPos t
  Child :: Int -> RPos child -> RPos (Node n children)
-- Approximation (GADTs omitted for readability):
data RPos = Here | GoChild Int RPos
\end{lstlisting}

\begin{fidelitynote}
In Lean, \lean{RPos} is indexed by the tree.  The \lean{child} constructor
is only well-typed when the tree is a \lean{node\,n\,children}, and $i$
must be a valid index in \lean{Fin\,n}.  Haskell's approximation loses
both the tree-index and the bound on $i$.
\end{fidelitynote}


%===================================================================
\section{The directed-container operations}
%===================================================================

% --- 3. RTree.sub -------------------------------------------------
\subsection{Definition: \lean{RTree.sub}}

\begin{defn}[Sub-tree at a position]
The sub-tree at position $p$ in tree $t$:
\begin{align*}
  t \downarrow \lean{here} &= t, \\
  (\lean{node}\,n\,\mathrm{children}) \downarrow \lean{child}(i, p)
    &= \mathrm{children}(i) \downarrow p.
\end{align*}
\end{defn}

\paragraph{Lean.}
\begin{lstlisting}[style=lean]
def RTree.sub : (t : RTree) -> RPos t -> RTree
  | t, .here => t
  | .node _ children, .child i p => (children i).sub p
\end{lstlisting}

\paragraph{Haskell.}
\begin{lstlisting}[style=haskell]
sub :: RTree -> RPos -> RTree
sub t             Here         = t
sub (Node _ cs)  (GoChild i p) = sub (cs !! i) p
sub (Node 0 _)   (GoChild _ _) = error "leaf has no children"
\end{lstlisting}


% --- 4. RTree.shift -----------------------------------------------
\subsection{Definition: \lean{RTree.shift}}

\begin{defn}[Shift (path concatenation)]
The shift $p \oplus q$ embeds position $q$ in the sub-tree at $p$ back
into the whole tree:
\begin{align*}
  t.\lean{shift}(\lean{here}, q) &= q, \\
  (\lean{node}\,n\,\mathrm{children}).\lean{shift}(\lean{child}(i, p), q)
    &= \lean{child}(i,\; \mathrm{children}(i).\lean{shift}(p, q)).
\end{align*}
\end{defn}

\paragraph{Lean.}
\begin{lstlisting}[style=lean]
def RTree.shift : (t : RTree) -> (p : RPos t) -> RPos (t.sub p) -> RPos t
  | _, .here, q => q
  | .node _ children, .child i p, q => .child i ((children i).shift p q)
\end{lstlisting}

\paragraph{Haskell.}
\begin{lstlisting}[style=haskell]
shift :: RTree -> RPos -> RPos -> RPos
shift _            Here          q = q
shift (Node _ cs) (GoChild i p) q = GoChild i (shift (cs !! i) p q)
\end{lstlisting}


%===================================================================
\section{Directed-container laws}
%===================================================================

% --- 5. RTree.shift_root (D3) ------------------------------------
\subsection{Theorem: \lean{RTree.shift\_root} (D3)}

\begin{thm}[D3: right identity]
Shifting by the root of the sub-tree recovers the original position:
\[
  t.\lean{shift}(p, \lean{here}) \;=\; p
  \quad\text{for all } p : \mathrm{RPos}(t).
\]
\end{thm}

\paragraph{Lean.}
\begin{lstlisting}[style=lean]
def RTree.shift_root : (t : RTree) -> (p : RPos t) ->
    t.shift p .here = p
\end{lstlisting}

\paragraph{Haskell.}
\begin{lstlisting}[style=haskell]
-- D3: prop_shift_root t p = shift t p Here == p
\end{lstlisting}


% --- 6. RTree.sub_shift (D4) -------------------------------------
\subsection{Theorem: \lean{RTree.sub\_shift} (D4)}

\begin{thm}[D4: sub-trees nest along the shift]
\[
  t \downarrow (p \oplus q) \;=\; (t \downarrow p) \downarrow q
  \quad\text{for all } p : \mathrm{RPos}(t),\; q : \mathrm{RPos}(t \downarrow p).
\]
\end{thm}

\paragraph{Lean.}
\begin{lstlisting}[style=lean]
def RTree.sub_shift : (t : RTree) -> (p : RPos t) -> (q : RPos (t.sub p)) ->
    t.sub (t.shift p q) = (t.sub p).sub q
\end{lstlisting}

\paragraph{Haskell.}
\begin{lstlisting}[style=haskell]
-- D4: prop_sub_shift t p q = sub t (shift t p q) == sub (sub t p) q
\end{lstlisting}


% --- 7. RTree.shift_assoc (D5) -----------------------------------
\subsection{Theorem: \lean{RTree.shift\_assoc} (D5)}

\begin{thm}[D5: shift is associative]
Path concatenation is associative (with transport along D4):
\[
  t.\lean{shift}(p \oplus q,\; D4^{-1} \triangleright r)
  \;=\;
  t.\lean{shift}(p,\; (t \downarrow p).\lean{shift}(q, r))
\]
for all $p : \mathrm{RPos}(t)$, $q : \mathrm{RPos}(t \downarrow p)$,
$r : \mathrm{RPos}((t \downarrow p) \downarrow q)$.
\end{thm}

\paragraph{Lean.}
\begin{lstlisting}[style=lean]
def RTree.shift_assoc : (t : RTree) -> (p : RPos t) -> (q : RPos (t.sub p)) ->
    (r : RPos ((t.sub p).sub q)) ->
    t.shift (t.shift p q) ((t.sub_shift p q).symm [> r) =
      t.shift p ((t.sub p).shift q r)
\end{lstlisting}

\paragraph{Haskell.}
\begin{lstlisting}[style=haskell]
-- D5 (transport suppressed):
-- prop_shift_assoc t p q r =
--   shift t (shift t p q) r == shift t p (shift (sub t p) q r)
\end{lstlisting}

\begin{fidelitynote}
As in \lean{ProofSearch}, the Lean statement requires $r$ to be transported
along $\mathrm{sub\_shift}$ (D4) before the left-hand shift can consume it.
In the rose-tree case, the \lean{child} constructor recurses into the
appropriate child and the recursion is structurally simpler than for
\lean{STree} (one constructor instead of \lean{left}/\lean{right}).
\end{fidelitynote}


%===================================================================
\section{The directed container}
%===================================================================

% --- 8. proofSearchN ----------------------------------------------
\subsection{Definition: \lean{proofSearchN}}

\begin{defn}[The \lean{proofSearchN} directed container]
Define the directed container of $N$-ary proof-search trees as:
\[
  \mathrm{proofSearchN} \;=\; (\mathrm{RTree},\; \mathrm{RPos},\;
  \lean{here},\; \lean{RTree.sub},\; \lean{RTree.shift})
\]
with:
\begin{itemize}
  \item D1: $t \downarrow \lean{here} = t$ (by $\mathrm{rfl}$);
  \item D2: $t.\lean{shift}(\lean{here}, q) = q$ (by $\mathrm{rfl}$);
  \item D3: $\lean{RTree.shift\_root}$;
  \item D4: $\lean{RTree.sub\_shift}$;
  \item D5: $\lean{RTree.shift\_assoc}$.
\end{itemize}
\end{defn}

\paragraph{Lean.}
\begin{lstlisting}[style=lean]
def proofSearchN : DirectedContainer where
  Shape := RTree
  Pos := RPos
  root := fun _ => .here
  sub := RTree.sub
  shift := RTree.shift
  d1 := fun _ => rfl
  d4 := RTree.sub_shift
  d2 := fun _ _ => rfl
  d3 := RTree.shift_root
  d5 := RTree.shift_assoc
\end{lstlisting}

\paragraph{Haskell.}
\begin{lstlisting}[style=haskell]
proofSearchN :: DirectedContainer RTree RPos
proofSearchN = DirectedContainer
  { shape = RTree
  , pos   = RPos
  , root  = const Here
  , sub   = sub
  , shift = shift
  -- D1: sub t Here = t              (by definition)
  -- D2: shift t Here q = q          (by definition)
  -- D3: shift t p Here = p          (RTree.shift_root)
  -- D4: sub t (shift t p q) = sub (sub t p) q   (RTree.sub_shift)
  -- D5: shift t (shift t p q) r = shift t p (shift (sub t p) q r)
  }
\end{lstlisting}


%===================================================================
\section{Translation conventions}
%===================================================================

\begin{enumerate}
\item \textbf{Fin-indexed children.}  Lean uses $\mathrm{Fin}(n) \to \mathrm{RTree}$
  for children, enabling structural recursion without well-founded recursion
  or \lean{List}.  Haskell uses a plain list; the branching-factor bound is
  not enforced.

\item \textbf{Structural recursion.}  Both \lean{RTree.sub} and \lean{RTree.shift}
  recurse on paths by structural induction, with \lean{here} as the base case
  and \lean{child i p} as the recursive case.  This is simpler than the
  binary case (one recursive constructor vs.\ two).

\item \textbf{Relationship to binary trees.}  \lean{proofSearch} is a
  special case of \lean{proofSearchN} with branching factors restricted to
  $\{0, 2\}$.  The Lean formalisation keeps them separate to avoid the
  round-trip gap that would arise from an embedding.

\item \textbf{D1 and D2 definitional.}  As in the binary case, the first
  case (\lean{here}) in each recursive definition is definitionally equal to
  the expected value, so D1 and D2 are proved by $\mathrm{rfl}$.

\item \textbf{Proof bodies omitted.}  Only statements are translated.
\end{enumerate}

\end{document}
